% =====================================================================
\chapter{Orientaciones de corrección del proyecto}
\label{ap:correccion}
% =====================================================================

Este apéndice está pensado para el profesor y para el propio estudiante, que debería aplicarlo a su trabajo antes de entregarlo. La corrección no busca una cifra de mejora concreta, sino comprobar que la cifra que se presenta es cierta, que se ha medido contra referencias honradas y que el sistema no viola nada que no pueda violar.

\section{Orden de la corrección}

\begin{enumerate}[label=\arabic*.]
  \item \emph{Reproducibilidad}. Se clona el repositorio en una máquina limpia y se ejecuta la orden única en el modo reducido documentado. Se comprueba que la huella del conjunto de prueba coincide con la registrada en la semana~1 y que las cifras principales del informe se reproducen con una diferencia relativa menor del 0{,}5\,\%. Si no se reproducen, el criterio «Código y demostración» se califica como insuficiente y las cifras no reproducidas no se tienen en cuenta en el resto de la rúbrica.
  \item \emph{Criterios obligatorios}. Se recorren los criterios A, B o C de la opción, con la evidencia que señala el informe (tablas, registros de MLflow, salidas de las pruebas automáticas).
  \item \emph{Rúbrica}. Se puntúa cada criterio con la tabla de la opción. En caso de duda entre dos niveles, decide la calidad de la justificación escrita, no el resultado numérico.
  \item \emph{Defensa}. Las preguntas de la demostración sirven para confirmar la autoría y detectar deuda técnica oculta~\cite{sculley}; también pueden mover la nota de cada criterio un nivel en cualquier sentido.
\end{enumerate}

\section{Comprobación de la evaluación honesta}

\begin{table}[H]
\centering\small
\caption{Cómo se comprueban las reglas R1 a R7 de la \cref{sec:proy-comun}.}\label{tab:honesta}
\begin{tabularx}{\textwidth}{@{}c>{\raggedright\arraybackslash}X>{\raggedright\arraybackslash}X@{}}
\toprule
Regla & Qué se mira & Señal de alarma \\
\midrule
R1 & El generador acepta la semilla y reproduce el escenario nominal; los parámetros coinciden con los del enunciado & Parámetros cambiados para facilitar la mejora sin justificarlo \\
R2 & Fecha y huella del conjunto de prueba; registros de experimentos & Ejecuciones sobre la prueba antes de la semana~7; ajuste de hiperparámetros con las métricas de prueba \\
R3 & Tabla con línea base, alternativa sencilla, optimizador y techo & Solo se compara con la práctica actual; techo ausente \\
R4 & Mismas semanas, trayectorias o días y mismas semillas en todas las políticas & Políticas evaluadas en periodos o con ruidos distintos \\
R5 & Intervalo $t$ pareado y \emph{bootstrap}; comprobación de autocorrelación & Desviación típica de cada política en lugar de la de la diferencia \\
R6 & Prueba automática que viola a propósito una restricción y comprueba que la salvaguarda actúa & Restricciones duras solo dentro del modelo de optimización \\
R7 & Declaración del uso de asistentes y de cómo se verificó lo generado & Código que el estudiante no sabe explicar en la defensa \\
\bottomrule
\end{tabularx}
\end{table}

\section{Errores frecuentes y su efecto en la nota}

\begin{table}[H]
\centering\small
\caption{Errores frecuentes en los proyectos y criterio afectado.}\label{tab:errores-proyecto}
\begin{tabularx}{\textwidth}{@{}>{\raggedright\arraybackslash}X>{\raggedright\arraybackslash}p{3.8cm}>{\raggedright\arraybackslash}p{3.0cm}@{}}
\toprule
Error & Remisión & Criterio afectado \\
\midrule
Big-M de $10^{6}$ «por si acaso»; \emph{gap} sin informar & M24, §2.3 y §2.4 & formulación \\
Previsión evaluada solo por su error, no por la decisión que induce & M24, §6.1 & componente de IA \\
Varianza predictiva en lugar de latente en la mejora esperada & M25, §4.5 & componente de IA \\
Episodios truncados tratados como terminales en el RL & M26, §4.4 & control \\
Comparar un MPC con previsiones con un PID sin ellas como si fuera una comparación de métodos & M26, §5.1 & evaluación \\
Restricción de rendimiento sobre la media del tiempo de fotograma & M27, §5.1 & selección con presupuesto \\
LOD sin validar: caras degeneradas, orientación incoherente & M27, §6.2 & cadena de LOD \\
Censura ignorada: vidas interrumpidas tratadas como fallos o descartadas & M28, §4.5 & componente de IA \\
Sesgo del ganador: informar la puntuación de validación del mejor candidato & M28, §2.2 & evaluación \\
VAN calculado frente a la práctica actual y atribuido entero a la IA & M28, §5.1 & valor \\
\bottomrule
\end{tabularx}
\end{table}

\section{Cifras de control}

Las siguientes magnitudes se deducen del escenario nominal con los resultados de los apuntes y permiten detectar generadores o gemelos mal construidos. No son resultados esperados del optimizador.
\begin{itemize}[label=\textendash]
  \item \emph{Opción (a)}. En régimen permanente, mantener el horno a 530~°C con 20~°C de ambiente exige $(530-20)/R=510$~kW, y durante la primera hora de una carga fría unos 150~kW más; con los dos centros de mecanizado en marcha (300~kW), un calentamiento a plena potencia (800~kW) supera la potencia contratada, de modo que esa restricción debe estar activa en algunas horas del plan. La constante de tiempo del horno es $RC=3$~h.
  \item \emph{Opción (b)}. Todas las texturas en BC7 y BC5 a resolución completa ocupan $40\cdot 3\cdot 4096^{2}\cdot\tfrac43$~bytes $\approx 2{,}68$~GB, por encima del presupuesto de 1{,}5~GB. Con $H=1080$ y $60^{\circ}$ de campo de visión, un error de 1~cm alcanza un píxel a $9{,}35$~m (M27, proposición~5.5). A 60~Hz el presupuesto es de $16{,}7$~ms (M27, definición~5.1).
  \item \emph{Opción (c)}. Con la demanda media, la capacidad exige unas $650\cdot 2{,}5/70=23{,}2$ furgonetas, y los lunes, $27{,}9$. Sin covariables, la política por edad de la proposición~4.13 de M28 con $\beta=2{,}5$, $\theta=15\,000$~km, $c_p=350$~€ y $c_f=1800$~€ tiene su óptimo en $T^{\ast}\approx 7300$~km, con un coste de unos 81{,}9~€ por cada 1000~km, frente a 87{,}3~€ con la revisión cada 10\,000~km y 135~€ sin mantenimiento preventivo. La probabilidad de fallar antes de 10\,000~km es $1-e^{-(2/3)^{2{,}5}}=0{,}30$.
\end{itemize}

\section{Preguntas de la defensa}

\begin{enumerate}[label=\arabic*.]
  \item \emph{Comunes}. ¿Qué parte de la mejora frente a la práctica actual se debe a la IA y qué parte a la alternativa sencilla? ¿Qué fracción del margen hasta el techo se captura y qué haría falta para capturar más? ¿Qué semana del conjunto de prueba es la peor y por qué? ¿Qué ocurre si se duplica el ruido del generador?
  \item \emph{Opción (a)}. ¿Cuál es el precio sombra de la potencia contratada y qué decisión de compra de potencia justificaría? ¿Por qué no basta planificar con el precio medio previsto? Muestre el filtro de seguridad actuando con un sensor congelado.
  \item \emph{Opción (b)}. ¿Qué cota da la relajación lagrangiana en un fotograma y qué hueco de dualidad tiene? ¿Cómo se ajustaron las ganancias del controlador de $\lambda$ y qué pasa si la GPU es un 30\,\% más lenta? ¿Qué comprobación geométrica falló más a menudo durante el desarrollo?
  \item \emph{Opción (c)}. ¿Cómo trata el modelo las vidas censuradas? ¿Cuánto cambia $T^{\ast}$ si $c_f$ se duplica? ¿Coincide la estimación por ponderación inversa de la propensión con el gemelo, y por qué podría no hacerlo?
\end{enumerate}
